\chapter{Quantification du risque, rentabilité et aide à la décision}

Ce chapitre exploite les modèles de probabilité de défaut et de perte en cas de défaut présentés au chapitre précédent pour construire les indicateurs qui éclairent effectivement la décision de crédit et le pilotage du portefeuille. Il détermine d'abord l'exposition au moment du défaut, puis combine ces trois composantes pour évaluer les pertes attendues selon IFRS 9, le capital économique et la rentabilité. Plusieurs analyses complémentaires enrichissent ensuite le pilotage du portefeuille, avant que le chapitre ne referme cette quantification sur le mécanisme de décision qui en fait usage à l'Origination, mécanisme dont le chapitre suivant présente l'architecture logicielle.

\section{Détermination de l'exposition au défaut (EAD)}

À la différence de la probabilité de défaut et de la perte en cas de défaut, présentées au chapitre précédent, l'exposition au moment du défaut (EAD) n'est pas estimée par un modèle statistique dans ce projet. Ce choix est délibéré : pour les crédits amortissables, l'EAD peut être déterminée à partir des caractéristiques contractuelles et de l'échéancier du crédit. Pour les lignes renouvelables, elle repose sur le taux d'utilisation lorsqu'il est disponible ou sur un facteur forfaitaire dans le cas contraire. Le projet privilégie ainsi des règles de calcul déterministes plutôt qu'un modèle statistique supplémentaire.

Pour les crédits amortissables (consommation, équipement, immobilier), l'EAD correspond au capital restant dû, calculé à partir de l'échéancier actuariel à annuités constantes établi lors de l'origination du crédit. Un crédit arrivé à son terme (ancienneté supérieure ou égale à sa durée totale) a ainsi une EAD nulle : son capital est intégralement remboursé, et il ne contribue plus à la perte attendue du portefeuille.

Pour les lignes de trésorerie renouvelables (BFR, propres au segment Entreprises), le calcul diffère selon la disponibilité du taux d'utilisation courant de la ligne. Lorsque ce taux est connu, l'EAD est obtenue en appliquant directement ce taux au montant de la ligne : $EAD = \text{Montant de la ligne} \times \text{Taux d'utilisation}$. Lorsqu'il ne l'est pas, l'EAD retient un facteur de conversion forfaitaire (CCF, \emph{Credit Conversion Factor}) de 0{,}75, valeur retenue dans le paramétrage du projet, appliqué à la ligne dans son ensemble : $EAD = \text{Montant de la ligne} \times 0{,}75$. Il convient de noter que cette seconde formule s'écarte de la convention réglementaire consistant à appliquer le CCF à la seule part non utilisée de la ligne, en s'ajoutant au montant déjà tiré : le projet retient ici une application plus simple du facteur de conversion, à la ligne entière plutôt qu'à son seul solde disponible.

\section{Évaluation des pertes attendues selon IFRS 9}

\subsection{Classification par stage et détection du SICR}

Dans l'implémentation du projet, le classement d'un compte dans l'un des trois stages IFRS 9, présentés au chapitre 1, repose sur la détection d'une augmentation significative du risque de crédit (SICR) depuis sa comptabilisation initiale. Cette détection est traduite en règles déterministes propres à chaque segment, inspirées des principes de la norme et adaptées aux données disponibles, dont le déclenchement d'un seul suffit à faire basculer un compte du Stage 1 au Stage 2 : pour le segment Particuliers, un retard de paiement supérieur ou égal à 30 jours, une proportion de mois en retard supérieure à 20~\% depuis l'origination du compte, un ratio de couverture par la garantie inférieur à 80~\%, ou une restructuration déjà intervenue ; pour le segment Entreprises, une dégradation de deux crans ou plus de la notation depuis l'origination, un \emph{Z"-score} d'Altman inférieur à 1,23 (zone de détresse), un choc macroéconomique défavorable, défini par la conjonction d'une hausse du chômage supérieure à 1~point et d'une croissance négative, ou un taux d'utilisation d'une ligne de trésorerie supérieur à 80~\%. Le Stage 3 est réservé aux comptes en défaut avéré. Lorsqu'une donnée nécessaire à un critère est indisponible, celui-ci n'est pas déclenché automatiquement. Cette règle évite de considérer une information manquante comme une preuve de dégradation, mais peut également limiter la détection de certaines situations de risque lorsque les informations nécessaires ne sont pas disponibles.

\subsection{Calcul différencié des pertes attendues en suivi de portefeuille}

Pour le suivi du portefeuille (\emph{Portfolio Analytics}), le calcul de la perte attendue applique effectivement la différenciation des horizons de calcul prévue par IFRS 9 selon le stage du compte. Au Stage 1, l'ECL retient la probabilité de défaut à douze mois, présentée au chapitre précédent. Au Stage 2, l'horizon de calcul est étendu à la durée de vie résiduelle de l'exposition : la probabilité de défaut correspondante est alors obtenue par un modèle de survie (régression de Cox), sous la forme d'une probabilité cumulée unique qu'un défaut survienne avant l'échéance résiduelle du crédit, plutôt que dans les douze mois suivants. Au Stage 3, le défaut étant déjà avéré, la probabilité de défaut n'est plus estimée mais fixée à 1 ; la perte en cas de défaut continue néanmoins d'être estimée par le modèle présenté au chapitre 3, plutôt que forfaitée à 100~\%, afin de refléter les récupérations encore attendues sur un compte déjà en défaut. Dans les trois cas, l'ECL retenue est le produit direct de cette probabilité de défaut, de la perte en cas de défaut et de l'exposition au moment du défaut ; il s'agit, pour les Stages 2 et 3, d'une estimation fondée sur une probabilité de défaut lifetime globale, et non d'un calcul décomposant la durée de vie résiduelle en pertes périodiques actualisées, ni d'une estimation des flux de trésorerie contractuels et recouvrables attendus telle que prévue par la norme pour les actifs dépréciés.

\subsection{Simplification effectivement appliquée à l'Origination}

Lors de l'évaluation d'une nouvelle demande de crédit, le stage du client n'est pas pour autant figé au Stage 1 par convention : il est recalculé à partir de son historique disponible, selon les mêmes critères SICR que ceux présentés ci-dessus. Pour un nouveau client ne disposant d'aucun historique exploitable, les critères fondés sur l'évolution du risque depuis l'origination ne peuvent pas être évalués. En l'absence d'un autre critère de classement, le compte est alors affecté au Stage 1.

En revanche, le calcul de la perte attendue à l'Origination ne reprend pas la différenciation par stage présentée pour le suivi de portefeuille : il retient systématiquement la probabilité de défaut à douze mois, quel que soit le stage du client. Le stage intervient néanmoins ailleurs dans la décision, à travers les règles présentées en section~4.5 (un Stage 3 entraînant un refus automatique, par exemple). Il s'agit là d'une simplification assumée du calcul de la perte attendue à l'Origination, et non d'une application intégrale du mécanisme de dépréciation par stage prévu par la norme, qui reste, pour sa part, réservé au suivi de portefeuille.

\section{Capital économique et rentabilité}

\subsection{Capital forfaitaire et RAROC à l'Origination}

La rentabilité d'un crédit, présentée au chapitre 1, s'apprécie au regard du capital qu'il immobilise : le RAROC (\emph{Risk-Adjusted Return on Capital}) rapporte le revenu net du crédit, diminué de la perte attendue et des frais de fonctionnement, au capital ainsi immobilisé. Pour chaque décision individuelle d'Origination, ce capital n'est pas calculé par simulation : le projet retient un forfait égal à 8~\% de l'exposition au moment du défaut. Cette valeur constitue une approximation opérationnelle utilisée pour le calcul du RAROC ; elle ne correspond pas à un calcul réglementaire des exigences de fonds propres selon Bâle II, dont le ratio minimal de 8~\% s'applique aux fonds propres rapportés aux actifs pondérés par les risques, non directement à l'EAD. Ce capital forfaitaire intervient également dans le calcul du taux cible proposé à l'utilisateur, qui intègre une rémunération de ce capital. Le taux plancher, en revanche, couvre uniquement le coût de financement, la perte attendue et les frais de fonctionnement, sans rémunération du capital immobilisé.

\subsection{Capital économique du portefeuille par simulation Monte Carlo}

Pour le pilotage du portefeuille, le capital économique est estimé différemment, par une simulation Monte Carlo de la distribution des pertes du portefeuille dans son ensemble. Les défauts individuels n'y sont pas simulés indépendamment : un modèle à facteur commun, inspiré des conventions de l'approche notations internes de Bâle II, introduit une corrélation entre les comptes d'un même segment, par l'intermédiaire d'un facteur macroéconomique partagé, avec une corrélation d'actifs de 0,12 pour les Particuliers et de 0,20 pour les Entreprises. De cette distribution simulée des pertes sont tirées la perte attendue (moyenne des simulations), la valeur en risque (VaR, perte au-delà de laquelle se trouvent 5~\% des scénarios simulés) et la perte inattendue (UL, \emph{Unexpected Loss}), définie comme l'écart entre la VaR et la perte attendue : $UL = VaR - EL$. Cette perte inattendue constitue le capital économique nécessaire pour absorber les pertes qui dépassent l'espérance.

Cette simulation appelle deux précisions. D'une part, son seuil de confiance (95~\%) n'est pas celui d'un calcul réglementaire de fonds propres, généralement fixé à 99,9~\% ; elle n'a donc pas vocation à se substituer à un tel calcul, mais à outiller le pilotage interne du portefeuille. D'autre part, l'horizon de la perte attendue simulée n'est pas unique : il hérite, pour chaque compte, de l'horizon de PD retenu par le calcul des pertes attendues présenté en section~4.2, douze mois en Stage 1, durée de vie résiduelle en Stage 2 et 3, ce qui en fait un horizon mixte plutôt qu'un horizon réglementaire explicite. Le capital ainsi obtenu par segment est ensuite réparti entre les comptes proportionnellement à leur exposition individuelle au sein de ce segment ; cette répartition est une allocation segmentaire simplifiée, et non une décomposition du risque de portefeuille attribuant à chaque compte sa contribution marginale exacte au risque global.

\subsection{Rentabilité : pricing et rôle du seuil de RAROC}

Pour chaque crédit, un taux plancher et un taux cible sont calculés à partir de la probabilité de défaut et de la perte en cas de défaut : le premier couvre les coûts du crédit sans rémunérer le capital qu'il immobilise, comme indiqué ci-dessus ; le second y ajoute la rémunération de ce capital au niveau du \emph{hurdle rate} du segment (13~\% pour les Particuliers, 15~\% pour les Entreprises), de sorte qu'un crédit tarifé à ce taux cible atteint exactement ce hurdle. Ce taux cible constitue la référence suggérée à l'utilisateur, rappelée lors de la décision présentée en section~4.5.

Dans la décision elle-même, le RAROC calculé pour le dossier n'est comparé qu'à un seuil minimal (\texttt{raroc\_min}), propre à chaque segment et à chaque palier de décision présenté en section~4.5 : c'est l'unique seuil de RAROC qui conditionne effectivement l'issue de la décision. Le hurdle rate, bien que calculé et conservé dans l'explication fournie à l'utilisateur pour situer le RAROC obtenu, n'intervient en revanche dans aucune règle de décision : sa comparaison au RAROC reste strictement informationnelle.

Il convient enfin de souligner que ces deux bases de capital, forfaitaire et issue de la simulation Monte Carlo, ne sont jamais combinées dans un même calcul de RAROC : le capital forfaitaire reste la seule base utilisée pour le RAROC d'une décision individuelle d'Origination, tandis que le capital Monte Carlo n'alimente que les indicateurs de rentabilité calculés à l'échelle du portefeuille.

\section{Analyses complémentaires du portefeuille}

\subsection{Segmentation du portefeuille par clustering}

La segmentation Particuliers/Entreprises, présentée au chapitre 3, sépare le portefeuille selon un critère économique de niveau d'encours, mais ne dit rien de la diversité des profils de risque à l'intérieur de chacun de ces deux segments. Une seconde segmentation, indépendante de la première, regroupe les comptes actifs d'un même segment selon la similarité de leur profil, afin d'identifier des sous-populations qui se comportent différemment face au risque.

Cette segmentation repose sur un algorithme de partitionnement non supervisé (K-means), appliqué séparément aux Particuliers et aux Entreprises, en quatre groupes. Les variables utilisées sont celles déjà construites pour la modélisation de la PD, présentée au chapitre 3, standardisées avant le partitionnement. Le choix de l'algorithme et du nombre de groupes a été établi empiriquement, en comparant K-means à un modèle de mélange gaussien selon deux dimensions complémentaires : la qualité de séparation des groupes, appréciée par le coefficient de silhouette, et la stabilité des partitions, évaluée par l'indice de Rand ajusté (\emph{Adjusted Rand Index}, ARI) entre deux initialisations différentes de l'algorithme. Le modèle de mélange gaussien produisait des groupes présentant une diversité comportementale plus importante, mais nettement moins stables (un indice de Rand ajusté descendant jusqu'à 0,39) ; K-means, à l'inverse, conserve une stabilité quasi parfaite (proche de 1) pour quatre groupes, un critère important pour limiter la dépendance des regroupements à l'initialisation de l'algorithme et obtenir des résultats reproductibles à données identiques.

Cette méthode a fait l'objet d'un ajustement a posteriori pour le segment Particuliers. Les variables macroéconomiques, communes à l'ensemble des comptes originés à une même période, dominaient initialement le partitionnement : deux des quatre groupes obtenus, représentant 77~\% du portefeuille Particuliers, reflétaient surtout le millésime d'origination des comptes plutôt qu'une réelle différence de profil de risque. Ces six variables macroéconomiques ont donc été retirées des variables utilisées pour ce segment uniquement. La stabilité du partitionnement en a été préservée, voire légèrement améliorée, et un groupe caractérisé par une délinquance comportementale demeure identifiable dans les deux configurations, avec un effectif et des caractéristiques proches. Dans la configuration retenue, il rassemble environ 1~150 comptes, avec une probabilité de défaut moyenne proche de 13~\% et un RAROC moyen proche de $-20~\%$, cette dernière valeur étant la moyenne arithmétique des RAROC individuels des comptes du groupe, calculés selon la méthode présentée en section~4.3, et non un RAROC recalculé sur les agrégats du groupe. Les quatre groupes obtenus sans ces variables macroéconomiques se sont par ailleurs révélés plus directement interprétables en termes de produit et de comportement de remboursement.

Pour le segment Entreprises, la même expérimentation retire également les variables macroéconomiques, avec une interprétabilité améliorée, mais une stabilité légèrement dégradée (indice de Rand ajusté de 0,95 contre 0,99 avec ces variables) ; aucun groupe de délinquance comportementale comparable à celui des Particuliers n'apparaît par ailleurs dans l'une ou l'autre configuration, les variables disponibles pour ce segment ne portant pas d'indicateur de retard comportemental équivalent à celles du segment Particuliers. En l'absence d'un gain aussi net que pour les Particuliers, la configuration initiale, incluant les variables macroéconomiques, a été conservée pour le segment Entreprises.

Les groupes ainsi obtenus doivent être interprétés avec prudence : ce sont des regroupements statistiques de comptes aux profils similaires, recalculés à chaque exécution plutôt que des catégories figées, et non des classes de risque réglementaires ou des catégories prédéfinies par la politique de crédit de la plateforme.

\subsection{Stress testing du portefeuille}

Au-delà du capital économique estimé par simulation Monte Carlo en section~4.3, qui quantifie une perte inattendue dans les conditions actuelles du portefeuille, un exercice de stress testing mesure l'effet d'une dégradation du contexte macroéconomique sur les mêmes indicateurs. Trois scénarios sont définis : un scénario de référence reprenant les conditions actuelles, un scénario adverse correspondant à un ralentissement économique modéré, et un scénario sévère correspondant à une crise systémique d'ampleur comparable aux chocs macroéconomiques les plus sévères traversés par la région ces dernières années. Chaque scénario est défini par un choc sur le chômage, la croissance, l'inflation et le taux directeur, présenté au tableau~\ref{tab:stress-macro}, et par un multiplicateur appliqué à la probabilité de défaut de chaque segment, présenté au tableau~\ref{tab:stress-pd}, plus marqué pour les Entreprises que pour les Particuliers. Ces chocs et ces multiplicateurs constituent des hypothèses de travail, choisies par jugement d'expert dans le cadre du projet, auxquelles le projet associe un niveau de confiance moyen, et non une calibration économétrique fondée sur des séries macroéconomiques historiques ni une valeur reprise telle quelle d'une publication identifiée.

\begin{table}[htbp]
    \centering
    \footnotesize
    \caption{Chocs macroéconomiques par scénario.}
    \label{tab:stress-macro}
    \begin{tabular}{p{2.6cm}p{2.6cm}p{2.6cm}p{2.6cm}p{2.6cm}}
        \toprule
        \textbf{Scénario} & \textbf{Chômage} & \textbf{Croissance} & \textbf{Inflation} & \textbf{Taux directeur} \\
        \midrule
        Référence & 0 pp & 0 pp & 0 pp & 0 pp \\
        Adverse & $+0,95$ pp & $-5,70$ pp & $+2$ pp & $+1$ pp \\
        Sévère & $+1,10$ pp & $-5,90$ pp & $+4$ pp & $+2$ pp \\
        \bottomrule
    \end{tabular}
\end{table}

\begin{table}[htbp]
    \centering
    \footnotesize
    \caption{Multiplicateur de la probabilité de défaut par scénario et par segment.}
    \label{tab:stress-pd}
    \begin{tabular}{p{3cm}p{3cm}p{3cm}}
        \toprule
        \textbf{Scénario} & \textbf{Particuliers} & \textbf{Entreprises} \\
        \midrule
        Référence & $\times 1$ & $\times 1$ \\
        Adverse & $\times 1,9$ & $\times 2,1$ \\
        Sévère & $\times 3,1$ & $\times 3,8$ \\
        \bottomrule
    \end{tabular}
\end{table}

Les composantes de la perte attendue ne sont pas toutes stressées de la même manière. La probabilité de défaut est directement multipliée par le facteur du scénario, puis bornée à l'intervalle $[0,1]$. La perte en cas de défaut n'est, elle, jamais multipliée directement : un abattement (\emph{haircut}) est appliqué au ratio de couverture par la garantie, et c'est ce ratio dégradé qui est soumis au modèle de perte en cas de défaut présenté au chapitre 3. Cette approche permet de conserver ce modèle comme mécanisme unique de calcul de cette composante, y compris dans les scénarios stressés. L'exposition au moment du défaut n'est stressée que pour les lignes de trésorerie renouvelables du segment Entreprises, par une dégradation de leur taux d'utilisation ; les crédits amortissables et les lignes d'investissement des Entreprises conservent leur exposition contractuelle, hors périmètre de cet exercice.

Le stage IFRS 9 de chaque compte n'est, enfin, recalculé sous stress que partiellement : pour les comptes Entreprises disposant des données nécessaires, les critères de dégradation liés au contexte macroéconomique et à l'utilisation des lignes de trésorerie, présentés en section~4.2, sont réévalués avec les valeurs stressées, à l'aide du même moteur de détection que celui utilisé en conditions normales ; les autres critères, ainsi que l'ensemble des critères du segment Particuliers, restent fixés à leur valeur observée, faute d'un modèle de projection comportementale sous stress pour ces critères. Un test réalisé sur des données réelles du portefeuille illustre l'ampleur de cet effet : le stage final de 450 comptes diffère de leur stage initial dans le scénario sévère. La perte attendue et le capital économique sont enfin recalculés sur le portefeuille stressé, en rejouant la simulation Monte Carlo présentée en section~4.3. Les résultats permettent ainsi de comparer, pour chaque scénario, l'évolution de la perte attendue totale, du ratio de couverture du portefeuille et du capital économique nécessaire pour absorber les pertes inattendues.

\subsection{Projection de l'évolution du portefeuille}

Le stress testing présenté ci-dessus mesure l'effet d'un choc hypothétique sur le portefeuille actuel ; la projection répond à une question différente : que deviendrait la répartition du portefeuille entre les stages IFRS 9, sur un horizon de douze mois, si la politique d'octroi et le contexte macroéconomique restaient inchangés. Il s'agit explicitement d'une projection sous hypothèse de stabilité, au fil de l'eau, et non d'une prévision économique au sens propre ni d'un engagement chiffré sur l'évolution réelle du portefeuille.

Cette projection repose sur une matrice de migration estimée empiriquement, à partir des snapshots disponibles dans le Feature Store présenté au chapitre 3, plutôt que déduite d'hypothèses théoriques sur la probabilité de défaut. Cette matrice décrit, pour chaque stage initial, la répartition des comptes entre les différents stages observés au snapshot suivant. Elle est obtenue en comparant le stage des mêmes comptes à deux anciennetés de snapshot différentes : six et douze mois pour les Particuliers, deux et quatre trimestres pour les Entreprises. Ces comptes et leurs trajectoires de stage restent ceux du périmètre synthétique du projet, présenté au chapitre 3 : cette matrice ne constitue donc pas une mesure de migration observée sur un historique réel de défauts bancaires. Cette matrice est ensuite appliquée de façon récursive, par deux pas successifs pour les deux segments, afin d'atteindre l'horizon de douze mois, ce qui suppose que les probabilités de transition observées sur l'intervalle de référence demeurent stables d'un pas à l'autre ; cette hypothèse est cohérente avec le scénario de stabilité retenu pour cette projection, mais ne permet pas de représenter une évolution progressive des comportements de remboursement au cours de l'horizon projeté. La perte attendue projetée à cet horizon reste une projection à l'échelle du portefeuille, et non compte par compte : elle multiplie l'effectif projeté de chaque stage par la probabilité de défaut, la perte en cas de défaut et l'exposition moyennes actuellement observées pour ce stage, plutôt que de recalculer individuellement chaque compte à douze mois.

L'incertitude de cette projection est quantifiée par une procédure de rééchantillonnage (bootstrap), appliquée au niveau des comptes plutôt qu'au résultat final : un échantillon de comptes de référence est tiré avec remise, la matrice de migration est recalculée sur cet échantillon, puis la projection et la perte attendue projetée sont recalculées à leur tour, l'ensemble de cette procédure étant répété 500 fois afin d'obtenir la distribution des résultats projetés. Les 5\textsuperscript{e} et 95\textsuperscript{e} centiles de cette distribution définissent les bornes de l'intervalle présenté, sans constituer un intervalle de confiance au sens statistique du terme. Cette incertitude reflète ainsi la précision avec laquelle la matrice de migration elle-même est estimée à partir des comptes disponibles, plutôt qu'une marge ajoutée après coup au résultat central. Les scénarios de stress présentés ci-dessus peuvent par ailleurs être appliqués à la trajectoire projetée, en réutilisant les mêmes multiplicateurs de probabilité de défaut, sans dupliquer une seconde logique de choc macroéconomique.

Cette projection reste volontairement limitée dans son périmètre : elle ne recourt pas à la simulation Monte Carlo du capital économique présentée en section~4.3, ne projette pas de scénario macroéconomique prospectif, et ne produit pas de valeur en risque à cet horizon. Seule l'évolution de la répartition par stage et de la perte attendue qui en découle, à politique et environnement inchangés, est couverte.

\subsection{Simulation et optimisation de la politique d'octroi}

Les trois analyses précédentes portent sur le portefeuille tel qu'il existe. Une dernière analyse complémentaire se place du point de vue inverse, celui d'une évolution de la politique d'octroi elle-même : que deviendraient le RAROC, la perte attendue, l'exposition et le profit du portefeuille si la banque refusait les X~\% de comptes présentant la probabilité de défaut la plus élevée. Il s'agit d'une simulation destinée à éclairer une discussion de politique de crédit, et non d'une prévision certaine des gains qui seraient effectivement réalisés : elle retient notamment l'hypothèse simplificatrice d'un taux d'intérêt et d'une marge brute inchangés après un tel durcissement, sans modèle de tarification dynamique qui ajusterait ces paramètres au nouveau profil du portefeuille résiduel.

Concrètement, les comptes du portefeuille sont triés par probabilité de défaut décroissante, puis une proportion croissante d'entre eux, les plus risqués, est retirée du portefeuille simulé. Pour chaque taux de refus testé, la perte attendue, l'exposition, le profit net et le RAROC sont recalculés sur le portefeuille résiduel, et comparés à ceux du portefeuille complet. Ce RAROC résiduel est un RAROC agrégé, calculé sur les montants cumulés du portefeuille résiduel, et non une moyenne des RAROC individuels des comptes qui le composent, à la différence de l'indicateur présenté pour un groupe en section~4.4.1.

Cette simulation a fait l'objet d'une correction méthodologique réelle. Sa formule de RAROC définissait auparavant son propre calcul, indépendant de celui présenté en section~4.3, et omettait notamment le coût de financement de l'exposition ; l'écart mesuré entre les deux formules atteignait environ 87,5 points de RAROC, un RAROC ainsi surestimé qui avait été exposé tel quel sur le tableau de bord de la politique d'octroi. La simulation a depuis été corrigée pour réutiliser directement la fonction de calcul du RAROC présentée en section~4.3, commune au traitement de la simulation de portefeuille et à celui de l'Origination. Cette correction garantit l'utilisation d'une même formule de calcul, notamment pour la prise en compte du coût de financement, sans pour autant imposer un résultat identique lorsque les bases de capital diffèrent : dans cette simulation de politique d'octroi, le capital retenu suit le paramétrage choisi, soit le forfait de 8~\% de l'exposition, soit une allocation du capital économique issu de la simulation Monte Carlo du portefeuille, tandis que l'Origination retient systématiquement le forfait, présenté en section~4.3.

\section{Mécanisme de décision — Origination}

\subsection{Vérification de l'éligibilité, préalable au calcul du risque}

Avant tout calcul quantitatif, une demande de crédit est d'abord soumise à une vérification d'éligibilité, portant sur le montant, la maturité demandée et, selon le segment, l'âge et le taux d'endettement pour les Particuliers, ou la concentration sectorielle pour les Entreprises, au-delà d'un seuil de 32~\% du portefeuille. Un refus à ce stade ne repose sur aucun modèle de risque, ni sur une probabilité de défaut calculée : il sanctionne le non-respect d'une règle de politique produit, indépendante du profil de risque du demandeur, et est tracé comme tel plutôt que présenté comme un refus motivé par le risque.

Les seuils effectivement appliqués diffèrent par segment. Pour les Particuliers, une demande est éligible si son montant ne dépasse pas 100~millions de XOF, sa maturité 120~mois, son taux d'endettement 50~\%, et si le demandeur est âgé de 18 à 70~ans. Pour les Entreprises, le montant maximal est porté à 5~milliards de XOF et la maturité maximale à 180~mois ; aucun seuil de taux d'endettement n'est en revanche appliqué à ce segment, ce ratio n'étant pas exploitable dans le périmètre actuel des données disponibles, exactement comme la concentration sectorielle ne s'applique pas aux Particuliers faute de secteur d'activité renseigné pour ce segment. Cette différence de traitement tient ainsi à la disponibilité et à la pertinence des indicateurs selon le segment, et non à une omission de la règle d'éligibilité. Ces seuils constituent des paramètres de politique produit propres au projet, et non une transposition de limites réglementaires ou de normes bancaires existantes.

\subsection{Calcul des indicateurs et application des règles de décision}

Pour les demandes éligibles, le stage IFRS 9 du client est d'abord déterminé selon la logique présentée en section~4.2. La plateforme calcule ensuite la probabilité de défaut, la perte en cas de défaut à partir du type de garantie renseigné pour le dossier, l'exposition au moment du défaut, la perte attendue et le RAROC, selon les principes présentés dans les sections précédentes.

La recommandation produite par la plateforme repose sur des règles hiérarchiques, appliquées dans un ordre déterminé. Un refus intervient en premier lieu si la probabilité de défaut dépasse un seuil maximal ou si le client est classé au Stage 3. Vient ensuite, pour le seul segment Entreprises, une règle de revue manuelle fondée sur ce même seuil de concentration sectorielle de 32~\% : en pratique, cette règle ne peut se déclencher que dans le cas limite d'une concentration égale à ce seuil, puisqu'une concentration qui le dépasse a déjà été écartée à l'étape d'éligibilité. Un accord sans condition est recommandé si la probabilité de défaut, le RAROC et le stage respectent simultanément des seuils favorables ; un accord sous conditions est recommandé pour des seuils intermédiaires sur ces mêmes critères, les conditions elles-mêmes étant dérivées des données réelles du dossier plutôt que d'une liste forfaitaire, par exemple l'exigence d'une garantie lorsque le dossier n'en comporte structurellement aucune. Les cas restants, ni clairement favorables ni clairement défavorables sur ces critères, sont orientés vers une revue manuelle plutôt que tranchés automatiquement.

Le tableau~\ref{tab:decision-policy} donne les valeurs effectivement appliquées par ces règles. Le seuil de RAROC n'y est pas unique par segment : il est plus exigeant pour un accord sans condition que pour un accord sous conditions, ce dernier tolérant un dossier moins rentable dès lors que des conditions viennent encadrer le risque résiduel. Comme pour les seuils d'éligibilité, ces valeurs constituent des paramètres de politique de crédit propres au projet, administrés dans une table de gouvernance plutôt que codés en dur, et non une calibration statistique ni une transposition de normes réglementaires.

\begin{table}[htbp]
    \centering
    \footnotesize
    \caption{Seuils de la politique de décision par segment et par palier.}
    \label{tab:decision-policy}
    \begin{tabular}{p{2.6cm}p{2.6cm}p{2.3cm}p{2.3cm}p{2.3cm}}
        \toprule
        \textbf{Palier} & \textbf{Segment} & \textbf{PD} & \textbf{RAROC min.} & \textbf{Stage max.} \\
        \midrule
        Accord & Particuliers & $\leq 5~\%$ & $\geq 10~\%$ & 1 \\
        Accord sous conditions & Particuliers & $\leq 12~\%$ & $\geq 5~\%$ & 2 \\
        Revue manuelle & Particuliers & $\leq 20~\%$ & — & — \\
        Refus & Particuliers & $> 20~\%$ & — & — \\
        \addlinespace
        Accord & Entreprises & $\leq 3~\%$ & $\geq 12~\%$ & 1 \\
        Accord sous conditions & Entreprises & $\leq 8~\%$ & $\geq 5~\%$ & 2 \\
        Revue manuelle & Entreprises & $\leq 15~\%$ & — & — \\
        Refus & Entreprises & $> 15~\%$ & — & — \\
        \bottomrule
    \end{tabular}
\end{table}

Ce tableau ne doit pas être lu indépendamment de la hiérarchie des règles présentée ci-dessus : le Stage 3 entraîne toujours un refus, y compris pour une probabilité de défaut par ailleurs faible, puisque cette règle est évaluée en premier. Un RAROC insuffisant, à l'inverse, n'entraîne jamais à lui seul un refus automatique : un dossier dont la probabilité de défaut reste sous le seuil de refus, mais dont le RAROC ne satisfait ni le seuil d'accord ni celui d'un accord sous conditions, est orienté vers une revue manuelle. Le RAROC constitue ainsi un critère de sélection entre les paliers d'accord, et non un motif autonome de refus automatique.

\subsection{Explicabilité, traçabilité et intervention humaine}

Chaque décision est accompagnée d'un rapport d'explication, qui rapproche la règle appliquée des indicateurs calculés et, lorsque le modèle de PD le permet, des facteurs SHAP présentés au chapitre précédent. L'impossibilité de produire une explication SHAP, par exemple si l'explication échoue techniquement, est elle-même tracée explicitement, plutôt que silencieusement ignorée. L'ensemble de la décision, ses indicateurs et son explication sont ensuite archivés dans le journal d'audit de la plateforme, accompagnés du rappel obligatoire que cette recommandation est générée automatiquement et que la décision finale relève de l'analyste habilité, conformément aux niveaux de délégation en vigueur.

Ce rappel prend une portée concrète lorsque la décision aboutit à une revue manuelle : un ticket est alors effectivement créé dans la file de travail des analystes, avec les indicateurs calculés pour le dossier. La plateforme ne tranche donc pas elle-même ces cas limites ; elle les transmet à un analyste habilité, auquel revient la décision finale.

\section*{Conclusion du chapitre}

Ce chapitre a présenté les mécanismes qui permettent de relier l'exposition au risque de crédit à l'évaluation des pertes attendues, du capital et de la rentabilité. La détermination de l'EAD complète ainsi les estimations de la probabilité de défaut et de la perte en cas de défaut pour permettre le calcul des pertes attendues. Ce calcul diffère toutefois selon le volet de la plateforme : le suivi de portefeuille applique des horizons différenciés selon le stage du compte, tandis que l'Origination retient systématiquement une probabilité de défaut à douze mois.

L'évaluation du capital suit également deux approches distinctes : un forfait opérationnel pour les décisions individuelles d'Origination et une estimation par simulation Monte Carlo pour le pilotage du portefeuille. Ces deux bases ne sont pas combinées dans un même calcul de rentabilité.

Au-delà de cette quantification, quatre analyses complémentaires enrichissent le pilotage du portefeuille, chacune répondant à une question distincte. La segmentation par clustering révèle une hétérogénéité des profils de risque que la seule distinction Particuliers/Entreprises ne permettait pas d'observer, jusqu'à isoler un groupe de comptes en délinquance comportementale. Le stress testing mesure la sensibilité du portefeuille à une dégradation hypothétique du contexte macroéconomique, quand la projection estime, à l'inverse, comment sa répartition par stage et sa perte attendue évolueraient sous l'hypothèse d'un contexte macroéconomique et d'une politique d'octroi inchangés. La simulation de politique d'octroi, enfin, relie ces analyses de portefeuille aux décisions de crédit individuelles, en estimant l'effet d'un durcissement de cette politique sur les mêmes indicateurs de rentabilité. Les analyses réalisées à l'échelle du portefeuille contribuent ainsi à son pilotage, tandis que les indicateurs calculés pour chaque dossier alimentent le mécanisme de décision d'Origination, qui formule une recommandation explicable et traçable sans se substituer à la décision finale de l'analyste habilité. Les simplifications retenues dans l'ensemble de ces traitements ont été précisées afin de délimiter la portée des résultats produits.

Le chapitre suivant présente l'architecture logicielle de la plateforme, l'organisation de ses composants et les interactions qui permettent de mettre en œuvre les mécanismes décrits dans les chapitres précédents.
